Sebagian besar situs web masih menjalankan aplikasi PHP pada versi 7.x atau lebih lama yang sudah tidak mendapatkan pembaruan keamanan resmi, sehingga sistem menghadapi akumulasi kerentanan yang terus bertambah tanpa kemungkinan diperbaiki. Migrasi manual pada ratusan \textit{file} kode tidak realistis, dan alat transformasi yang tersedia saat ini tidak menyertakan pemeriksaan keamanan sebagai bagian terintegrasi dari alur migrasi. Organisasi yang beroperasi di bawah standar ISO/IEC 27001:2022 juga harus melakukan audit keamanan sebagai langkah terpisah setelah migrasi selesai, sehingga menambah beban operasional.

\hspace{1cm} Penelitian ini membangun \textit{platform} berbasis Python yang mengintegrasikan konversi otomatis aplikasi web PHP lama ke versi PHP target yang dapat dikonfigurasi, dipadukan dengan pemeriksaan kepatuhan \textit{secure coding} berbasis ISO/IEC 27001:2022 dalam satu alur kerja \textit{open source} yang dapat dijalankan secara lokal. \textit{Platform} ini menggunakan Rector untuk transformasi kode berbasis AST, Semgrep untuk pemindaian keamanan, PHPStan untuk validasi pasca-konversi, serta lima model LLM \textit{open source} yang dijalankan secara lokal melalui Ollama---DeepSeek Coder 6.7b, Qwen2.5-Coder 7b, CodeLlama 7b, Mistral 7b, dan Llama 3.1 8b---untuk analisis semantik dan saran perbaikan kerentanan. Studi kasus dilakukan pada kode PHP milik FT UGM (Dashboard TCK, 245 \textit{file}) dan CBT DTI UGM (326 \textit{file}).

\hspace{1cm} Rector mencatat tingkat konversi 97,6\% pada dataset FT UGM dan 91,7\% pada dataset CBT DTI UGM, dengan tingkat validitas sintaks pasca-konversi 100\% pada kedua studi kasus. Nilai $\Delta F = 0$ mengonfirmasi bahwa proses konversi tidak memperkenalkan kerentanan baru. Laporan JSON yang dihasilkan secara otomatis mencakup pemetaan per temuan ke kontrol ISO/IEC 27001:2022 Annex A yang relevan dan dapat langsung digunakan sebagai dokumentasi audit. Dari eksperimen komparasi lima model LLM, Qwen2.5-Coder 7b direkomendasikan untuk integrasi ke dalam sistem otomatis karena menghasilkan \textit{Syntax Validity Rate} (SVR) 100\% secara konsisten di semua kondisi eksperimen. \textit{Platform} ini membuktikan bahwa integrasi antara konversi kode otomatis dan pemeriksaan kepatuhan ISO/IEC 27001:2022 dapat dijalankan sepenuhnya secara lokal tanpa mengirimkan kode sumber ke layanan eksternal.

\noindent{Kata kunci} : migrasi PHP, \textit{secure coding}, ISO/IEC 27001:2022, \textit{large language model}, transformasi kode AST
